\ifreview
\title[]{\MSAGA{}: Fusing \FA{} \\ within a Single Systolic Array}
\else
\title[]{\SA{}: Fusing \FA{} within a Single Systolic Array}
\fi


\author{Jiawei Lin}
\email{jiawei.lin@epfl.ch}
\affiliation{%
  \institution{EPFL}
  \city{Lausanne}
  \country{Switzerland}
}

\author{Yuanlong Li}
\email{yuanlong.li@epfl.ch}
\affiliation{%
  \institution{EPFL}
  \city{Lausanne}
  \country{Switzerland}
}

\author{Guokai Chen}
\email{guokai.chen@epfl.ch}
\affiliation{%
  \institution{EPFL}
  \city{Lausanne}
  \country{Switzerland}
}

\author{Thomas Bourgeat}
\email{thomas.bourgeat@epfl.ch}
\affiliation{%
  \institution{EPFL}
  \city{Lausanne}
  \country{Switzerland}
}

\renewcommand{\shortauthors}{Jiawei et al.}

\begin{abstract}

Transformer models rely heavily on the scaled dot-product attention (SDPA) operation,
typically implemented as \FA{}.
Characterized by its frequent interleaving of matrix multiplications and softmax operations,
\FA{} fails to fully utilize the compute resources of modern systolic-array-based accelerators
designed for consecutive and large matrix multiplications.

To fully unleash the performance potential of systolic arrays for \FA{},
we propose \MSAGA{}, an enhanced systolic array architecture that runs the entire \FA{}
on the array without external vector units.
Combined with \SA{}, an optimized kernel for \MSAGA{} that achieves fine-grained and element-wise
overlapping of \FA{} operations, \MSAGA{} maximizes array utilization while preserving the
original floating-point operation order of \FA{}.
We implement \MSAGA{} in synthesizable RTL and evaluate its performance against state-of-the-art
systolic-array-based accelerators.
Our results show that \MSAGA{} achieves 1.77$\times$ and 4.83$\times$ higher attention \FLOPs{} utilization
compared to AWS Neuron-v2 and Google TPUv5e, respectively.
We synthesize \MSAGA{} in a 16\,nm technology at 1.5\,GHz, and results indicate
only a 12\% area overhead compared to a standard weight-stationary systolic array.
\end{abstract}








\maketitle
